\documentclass[11pt]{article}
\usepackage[margin=1in]{geometry}
\usepackage[T1]{fontenc}
\usepackage[utf8]{inputenc}
\usepackage{amsmath,amssymb,mathtools}
\usepackage{booktabs}
\usepackage{hyperref}
\usepackage{xcolor}
\usepackage{listings}

\hypersetup{
  colorlinks=true,
  linkcolor=blue!60!black,
  urlcolor=blue!60!black
}

\definecolor{codebg}{RGB}{248,248,248}
\definecolor{codeborder}{RGB}{210,210,210}
\definecolor{codecomment}{RGB}{80,120,80}
\definecolor{codekeyword}{RGB}{30,60,150}

\lstdefinestyle{reportcode}{
  backgroundcolor=\color{codebg},
  basicstyle=\ttfamily\small,
  breaklines=true,
  columns=fullflexible,
  commentstyle=\color{codecomment},
  frame=single,
  framerule=0.5pt,
  keywordstyle=\color{codekeyword},
  rulecolor=\color{codeborder},
  keepspaces=true,
  showstringspaces=false,
  tabsize=2
}

\lstset{style=reportcode, language=Python}

\setlength{\parindent}{0pt}
\setlength{\parskip}{0.5em}

\newcommand{\R}{\mathbb{R}}
\newcommand{\ysp}{y_{\mathrm{sp}}}
\newcommand{\uvec}{u}
\newcommand{\xvec}{x}
\newcommand{\Aaug}{A_{\mathrm{aug}}}
\newcommand{\Baug}{B_{\mathrm{aug}}}
\newcommand{\Caug}{C_{\mathrm{aug}}}
\newcommand{\Qout}{Q_{\mathrm{out}}}
\newcommand{\Rin}{R_{\mathrm{in}}}
\newcommand{\Hp}{H_p}
\newcommand{\Hc}{H_c}
\newcommand{\argmin}{\mathop{\mathrm{arg\,min}}}


\title{Model Multipliers And Observer-Pole Adaptation}
\date{2026-04-01}

\begin{document}
\maketitle

\section*{Method Goal}

This report groups the notebook families that adapt the internal prediction
model rather than only the MPC cost. The main body covers TD3 and SAC methods
that scale the augmented state-space matrices. The final section covers the
observer-pole method, which keeps the plant model fixed but adapts the observer
gain through pole placement.

\section*{Notebook Sources Covered}

\begin{itemize}
\item \texttt{RL\_assisted\_MPC\_matrices.ipynb}
\item \texttt{RL\_assisted\_MPC\_matrices\_dist.ipynb}
\item \texttt{RL\_assisted\_MPC\_matrices\_SAC.ipynb}
\item \texttt{RL\_assisted\_MPC\_matrices\_dsit\_SAC.ipynb}
\item \texttt{RL\_assisted\_MPC\_matrices\_model\_mismatch.ipynb}
\item \texttt{RL\_assisted\_MPC\_Poles.ipynb}
\end{itemize}

\section*{Matrix-Multiplier Method}

\subsection*{State, Action, Reward}

The TD3 and SAC matrix notebooks reuse the standard scaled RL state
\[
s_t =
\begin{bmatrix}
\tilde{x}_{t} \\
\tilde{y}_{\mathrm{sp},t} \\
\tilde{u}_{t}^{\mathrm{dev}}
\end{bmatrix}.
\]

The action is continuous with dimension $1 + n_u$:
\[
a_t \mapsto [\alpha_t, \delta_{t,1}, \dots, \delta_{t,n_u}].
\]
After remapping from $[-1, 1]$, the notebook uses these values to scale the
augmented model matrices relative to immutable baselines:
\[
A_t^{\mathrm{now}} =
\begin{bmatrix}
\alpha_t A_{\mathrm{phys}} & * \\
* & *
\end{bmatrix},
\qquad
B_t^{\mathrm{now}} = B_{\mathrm{base}} \operatorname{diag}(\delta_t)
\]
in the implementation sense used by the notebooks:
\[
A_{\mathrm{now}}[:n_{\mathrm{phys}}, :n_{\mathrm{phys}}] *= \alpha_t,
\qquad
B_{\mathrm{now}}[:n_{\mathrm{phys}}, :] *= \delta_t.
\]

\subsection*{Notebook-local matrix update logic}

\begin{lstlisting}
m_model = map_to_bounds(a_model_raw, model_low, model_high)
alpha = float(m_model[0])
delta = np.asarray(m_model[1:1 + n_inputs], float).reshape(-1)

A_now = A_base.copy()
B_now = B_base.copy()
n_phys = n_states - n_outputs
A_now[:n_phys, :n_phys] *= alpha
B_now[:n_phys, :] *= delta.reshape(1, -1)

MPC_obj.A = A_now
MPC_obj.B = B_now
\end{lstlisting}

\subsection*{TD3 Versus SAC}

The TD3 notebooks use a deterministic actor with exploration noise and delayed
policy updates. The SAC notebooks keep the same high-level control-loop shape
but replace the actor with a Gaussian policy and entropy-regularized critic
targets.

\begin{lstlisting}
next_action, logp_next, mean_next = self.actor.sample(ns)
q_next = self.critic_target.combined_forward(ns, next_action, mode="min")
alpha = self.log_alpha.exp()
y = r + self.gamma * (1.0 - done) * (q_next - alpha * logp_next)
\end{lstlisting}

\section*{Mismatch Variant}

The mismatch notebook keeps the same matrix-scaling mechanism but changes the
operating conditions more aggressively. The purpose is to test whether online
adaptation of $(A, B)$ can keep the linear MPC layer useful when the plant moves
away from the nominal identification point.

\section*{Observer-Pole Adaptation}

\subsection*{Method Role}

The pole notebook does not scale $(A, B)$ directly. Instead, it lets a TD3
policy propose observer poles, sanitizes them, smooths them, and recomputes the
observer gain $L$ through pole placement.

\subsection*{Reusable observer-gain helper}

\begin{lstlisting}
def compute_observer_gain(A, C, desired_poles):
    obs_gain_calc = signal.place_poles(A.T, C.T, desired_poles,
                                       method='KNV0')
    L = np.squeeze(obs_gain_calc.gain_matrix).T
    observability_matrix = control.obsv(A, C)
    ...
    return L
\end{lstlisting}

\subsection*{Notebook-local pole update logic}

\begin{lstlisting}
if (i % pole_update_every) == 0:
    poles_prop = map_to_bounds(action, LOW_POLES, HIGH_POLES)
    poles_prop = sanitize_poles(poles_prop, LOW_POLES, HIGH_POLES)
    poles_cur = (1.0 - pole_smooth) * poles_cur + pole_smooth * poles_prop
    poles_cur = sanitize_poles(poles_cur, LOW_POLES, HIGH_POLES)
    try:
        L_new = compute_observer_gain(MPC_obj.A, MPC_obj.C, poles_cur)
        if np.all(np.isfinite(L_new)):
            L = L_new
    except Exception:
        pass
\end{lstlisting}

\section*{Core Algorithm Flow}

\begin{enumerate}
\item Build the scaled RL state from the observer state and operating point.
\item Query either a TD3 or SAC agent for model multipliers, or a TD3 agent for
      observer poles.
\item Remap the action to physical coefficient bounds.
\item Update $(A, B)$ or recompute $L$.
\item Solve MPC, apply the first move, step the plant, and roll the observer.
\item Store the transition and train the active RL agent after warm start.
\end{enumerate}

\section*{Limitations, Caveats, And Open Questions}

\begin{itemize}
\item The matrix and pole methods live mainly in notebook supervisors, not in a
      shared experiment module.
\item The SAC notebooks change the learning rule but not the surrounding control
      loop, so careful comparison depends on keeping all notebook-level settings
      aligned.
\item The repo also contains parallel TD3 implementations such as
      \texttt{agent\_modified.py}; notebook imports must be checked before
      extending this family.
\item Observer adaptation and matrix adaptation both depend strongly on the
      observer state, so stale helper modules should not be treated as the main
      execution path.
\end{itemize}

\end{document}
